% ======================================================================
% INF5163 – Projet de session 2 – Présentation Beamer (12 min + 7-8 min de questions)
% Compilation : Tectonic (LaTeX Workshop, à l'enregistrement) → presentation/main.pdf
% En ligne de commande, depuis presentation/ : tectonic -X compile main.tex
% Environ 14 diapos de contenu ; chaque membre présente le bloc de sa section.
% Grille : Beamer – structure et clarté (15) ; discussion critique (15) ;
%          temps et participation (10).
% ======================================================================
\documentclass[aspectratio=169,11pt]{beamer}

% --- Thème sobre ---------------------------------------------------------------
\usetheme{SimpleDarkBlue}
\usecolortheme{default}
\usefonttheme{professionalfonts}
\setbeamertemplate{navigation symbols}{}
\setbeamertemplate{footline}{%
  \hfill\usebeamerfont{footline}\usebeamercolor[fg]{footline}%
  \insertshortauthor\quad\insertframenumber{} / \inserttotalframenumber\hspace{1.5em}\vspace{0.8em}}
\setbeamercolor{footline}{fg=gray}
\setbeamertemplate{itemize items}[circle]
\setbeamertemplate{bibliography item}[text]

% --- Langue et encodage ---------------------------------------------------------
\usepackage[T1]{fontenc}
\usepackage{lmodern}
\usepackage[french]{babel}
\usepackage[autostyle=true]{csquotes}
\usepackage{booktabs}
\graphicspath{{../shared/figures/}}

% --- Macros partagées avec le rapport (bibliographie : ../shared/references.bib) ---------
% ======================================================================
% Macros partagées par le rapport, l'entente autonome et la présentation.
% Ce fichier est chargé (\input) dans les préambules : ne rien y composer.
% ======================================================================

% --- Noms propres et cadres ------------------------------------------------
% Utilisation : \fedesante{}, \pcap{}, \eptc{} (les accolades évitent
% que l'espace suivante soit avalée).
\newcommand{\fedesante}{FédéSanté}
% PCAP est une marque déposée du Centre de gouvernance de l'information des
% Premières Nations (CGIPN) ; le symbole ® est demandé par le CGIPN.
% TODO : vérifier la forme d'usage recommandée par le CGIPN.
\newcommand{\pcap}{PCAP\textsuperscript{\textregistered}}
% TODO : vérifier la version de l'EPTC 2 en vigueur à la date de remise
% et ajuster l'entrée correspondante dans shared/references.bib.
\newcommand{\eptc}{EPTC~2}
\newcommand{\loivingtcinq}{Loi~25}

% --- Espace réservé visible ----------------------------------------------
% Marque dans le PDF ce qui reste à rédiger. Aucun \arediger ne doit
% subsister dans la version remise (vérifier avec : grep -rn arediger).
\newcommand{\arediger}[1]{{\color{gray}\itshape[À rédiger : #1]}}

% --- Sigles et acronymes (paquet acronym) ----------------------------------
% \ac{FL} : forme longue + sigle au premier emploi, sigle ensuite.
% \acs{FL} : sigle seul ; \acl{FL} : forme longue seule.
\usepackage{acronym}
\acrodef{FL}[FL]{apprentissage fédéré (\emph{federated learning})}
\acrodef{DP}[DP]{confidentialité différentielle (\emph{differential privacy})}
\acrodef{FedAvg}[FedAvg]{moyenne fédérée (\emph{Federated Averaging})}
\acrodef{CER}[CER]{comité d'éthique de la recherche}
\acrodef{PCAP}[PCAP]{propriété, contrôle, accès et possession}


\title{Quand le modèle voyage}
\subtitle{Éthique, propriété et souveraineté dans un consortium d'apprentissage fédéré en santé}
% Membres en ordre alphabétique de nom de famille.
\author[Équipe FédéSanté]{\small Alex D'Andrea, Hamidou Dia, Gaston Raphaël Kapawa\\Ninette Mahoro, Moïse Armand Massack, Charlin Jordan Moukam Djongoue\\Frank Onana Tsanga, Stive Roussel Tankio Djiokap}
\institute{INF5163 -- Méthodologie de recherche en informatique\\UQO -- Prof.~N'dah Daniel Yapi, Ph.D.}
\date{Automne 2026}

\begin{document}

\begin{frame}[plain]
  \titlepage
\end{frame}

\begin{frame}{Plan}
  \tableofcontents
\end{frame}

% Un bloc par section du rapport ; chaque responsable présente sa section
% (présentateur·trice et durée en tête de chaque fichier de diapos/).
% !TEX root = ../main.tex
% ======================================================================
% Bloc : Introduction
% Présentateur·trice : Ninette Mahoro | Durée cible : 1 min
% (Total visé : 12 min ; respect du temps et participation = 10 pts.)
% ======================================================================
\section{Introduction}

\begin{frame}{Mise en situation : FédéSanté}
  % TODO : 12 établissements dont 3 centres de santé de Premières Nations ; prédiction des complications du diabète ; justification du sujet.
  \arediger{Mise en situation : FédéSanté}
\end{frame}

\begin{frame}{« Les données ne quittent jamais l'établissement »}
  % TODO : la promesse et la question critique qui guide la présentation.
  \arediger{« Les données ne quittent jamais l'établissement »}
\end{frame}


% !TEX root = ../main.tex
% ======================================================================
% Bloc : Problématique et objectifs
% Présentateur·trice : Hamidou Dia | Durée cible : 1 min
% (Total visé : 12 min ; respect du temps et participation = 10 pts.)
% ======================================================================
\section{Problématique et objectifs}

\begin{frame}{Problématique et objectifs}
  % TODO : ce que FédéSanté visait à résoudre ; objectifs SMART ; mandat de l'analyse.
  \arediger{Problématique et objectifs}
\end{frame}


% !TEX root = ../main.tex
% ======================================================================
% Bloc : Revue de littérature
% Présentateur·trice : Gaston Raphaël Kapawa | Durée cible : 1,5 min
% (Total visé : 12 min ; respect du temps et participation = 10 pts.)
% ======================================================================
\section{Revue de littérature}

\begin{frame}{Ce que dit la littérature}
  % TODO : apprentissage fédéré en santé, attaques et DP, gouvernance autochtone ; lacunes.
  \arediger{Ce que dit la littérature}
\end{frame}


% !TEX root = ../main.tex
% ======================================================================
% Bloc : Méthodologie
% Présentateur·trice : Frank Onana Tsanga | Durée cible : 1,5 min
% (Total visé : 12 min ; respect du temps et participation = 10 pts.)
% ======================================================================
\section{Méthodologie}

\begin{frame}{Architecture fédérée et confidentialité différentielle}
  % TODO : FedAvg pondéré par la taille des sites ; paramètres ε, δ, écrêtage, bruit.
  \arediger{Architecture fédérée et confidentialité différentielle}
\end{frame}


% !TEX root = ../main.tex
% ======================================================================
% Bloc : Intégrité, éthique et ÉDI
% Présentateur·trice : Moïse Armand Massack | Durée cible : 2 min
% (Total visé : 12 min ; respect du temps et participation = 10 pts.)
% ======================================================================
\section{Intégrité, éthique et ÉDI}

\begin{frame}{Ce que la promesse ne couvre pas}
  % TODO : inférence d'appartenance et inversion de modèle à partir des mises à jour.
  \arediger{Ce que la promesse ne couvre pas}
\end{frame}

\begin{frame}{Consentement et usages secondaires}
  % TODO : large, spécifique ou dynamique ? Protocole de gouvernance (EPTC 2, Loi 25).
  \arediger{Consentement et usages secondaires}
\end{frame}

\begin{frame}{ÉDI : petits sites, sous-groupes, stigmatisation}
  % TODO : agrégation pondérée par la taille, performance par sous-groupe, communautés plus touchées.
  \arediger{ÉDI : petits sites, sous-groupes, stigmatisation}
\end{frame}


% !TEX root = ../main.tex
% ======================================================================
% Bloc : Gestion, responsabilité et PI
% Présentateur·trice : Alex D'Andrea | Durée cible : 1,5 min
% (Total visé : 12 min ; respect du temps et participation = 10 pts.)
% ======================================================================
\section{Gestion, responsabilité et PI}

\begin{frame}{Propriété du modèle et PCAP}
  % TODO : à qui appartient le modèle global ? PCAP appliqué aux paramètres ; brevet ou code ouvert.
  \arediger{Propriété du modèle et PCAP}
\end{frame}

\begin{frame}{Retrait, désapprentissage et documentation}
  % TODO : un site peut-il exiger le désapprentissage ? Fiche de modèle, fiches de données, registre, accès externe.
  \arediger{Retrait, désapprentissage et documentation}
\end{frame}


% !TEX root = ../main.tex
% ======================================================================
% Bloc : Valorisation et diffusion
% Présentateur·trice : Stive Roussel Tankio Djiokap | Durée cible : 1,5 min
% (Total visé : 12 min ; respect du temps et participation = 10 pts.)
% ======================================================================
\section{Valorisation et diffusion}

\begin{frame}{Plateforme publique ou licence commerciale ?}
  % TODO : tension entre les deux voies ; partage des bénéfices avec les communautés.
  \arediger{Plateforme publique ou licence commerciale ?}
\end{frame}

\begin{frame}{Diffuser auprès des CER, des instances autochtones et des patients}
  % TODO : deux publics, deux messages ; volet vulgarisation.
  \arediger{Diffuser auprès des CER, des instances autochtones et des patients}
\end{frame}


% !TEX root = ../main.tex
% ======================================================================
% Bloc : Entente de gouvernance
% Présentateur·trice : Charlin Jordan Moukam Djongoue | Durée cible : 1,5 min
% (Total visé : 12 min ; respect du temps et participation = 10 pts.)
% ======================================================================
\section{Entente de gouvernance}

\begin{frame}{Entente de gouvernance : articles clés}
  % TODO : droits PCAP, retrait et désapprentissage, approbation des publications, partage des bénéfices, différends.
  \arediger{Entente de gouvernance : articles clés}
\end{frame}


% !TEX root = ../main.tex
% ======================================================================
% Bloc : Conclusion
% Présentateur·trice : ___ (non attribuée : à décider en équipe) | Durée cible : 0,5 min
% (Total visé : 12 min ; respect du temps et participation = 10 pts.)
% ======================================================================
\section{Conclusion}

\begin{frame}{La promesse suffit-elle ?}
  % TODO : réponse au mandat, limites et risques résiduels.
  \arediger{La promesse suffit-elle ?}
\end{frame}



% TODO : retirer ce \nocite quand les sources seront citées dans les diapos.
\nocite{eptc2,loi25,pcap,mcmahan2017fedavg}
\begin{frame}[allowframebreaks]{Références}
  \footnotesize
  \bibliographystyle{IEEEtran}
  \bibliography{../shared/references}
\end{frame}

\begin{frame}[plain]
  \centering
  \vfill
  {\Huge Questions ?\par}
  \vspace{1em}
  {\large Merci de votre attention\par}
  \vfill
\end{frame}

\end{document}
